\documentclass[12pt,a4paper]{article}

\usepackage[utf8]{inputenc}
\usepackage[T1]{fontenc}
\usepackage[english]{babel}
\usepackage{amsmath}
\usepackage{amsfonts}
\usepackage{amssymb}
\usepackage{graphicx}
\usepackage{listings}
\usepackage{xcolor}
\usepackage{hyperref}
\usepackage{geometry}

\geometry{margin=1in}

\definecolor{codegreen}{rgb}{0,0.6,0}
\definecolor{codegray}{rgb}{0.5,0.5,0.5}
\definecolor{codepurple}{rgb}{0.58,0,0.82}
\definecolor{backcolour}{rgb}{0.95,0.95,0.92}

\lstdefinestyle{mystyle}{
    backgroundcolor=\color{backcolour},
    commentstyle=\color{codegreen},
    keywordstyle=\color{magenta},
    numberstyle=\tiny\color{codegray},
    stringstyle=\color{codepurple},
    basicstyle=\ttfamily\footnotesize,
    breakatwhitespace=false,
    breaklines=true,
    captionpos=b,
    keepspaces=true,
    numbers=left,
    numbersep=5pt,
    showspaces=false,
    showstringspaces=false,
    showtabs=false,
    tabsize=2
}

\lstset{style=mystyle}

\title{
    \textbf{Implementation of a Minimax Agent for Mancuna} \\
    \large Computer Games \\
    Master in Computer Engineering\\
    \vspace{1cm}
    \includegraphics[width=0.4\textwidth]{ualg_logo.png} \\ % Note: Placeholder for logo
    \vspace{1cm}
}

\author{
    Alexandre Boavida \\
    Jan-Erik Voigt \\
    Philippe Schlup
}

\date{April 21, 2026}

\begin{document}

\maketitle

\begin{center}
    \textbf{University of Algarve} \\
    Faculty of Science and Technology \\
    \vspace{0.5cm}
    \textbf{Professor:} Helder Anicedo Amadeu de Sousa Daniel
\end{center}

\newpage

\tableofcontents

\newpage

\section{Introduction}
This report details the implementation of an artificial intelligence agent for the game of \textit{Mancuna}, a simplified variant of the traditional Mancala game played on a $2 \times 2$ board. The project was developed as part of the Computer Games course in the Master in Computer Engineering at the University of Algarve.

The objective was to create a competitive player capable of searching the game tree efficiently to make optimal moves. Our implementation, found in \texttt{PlayerMancuna.cs}, utilizes several advanced game-playing techniques, including Minimax search with Alpha-Beta pruning, a transposition table, and a carefully designed evaluation function that addresses the horizon effect.

\section{The Mancuna Game}
Mancuna is played on a board with two rows and two columns. Each player controls one row (two pits). In the standard configuration, each pit starts with 2 tokens, for a total of 8 tokens on the board.
\begin{itemize}
    \item \textbf{Movement:} A player selects a non-empty pit on their side and redistributes the tokens counter-clockwise.
    \item \textbf{Winning Condition:} A player wins if the opponent's side is empty and it's the opponent's turn to move.
    \item \textbf{Scoring:} The score is the total number of tokens on a player's side of the board.
\end{itemize}

\section{Implementation Details}

\subsection{Minimax Algorithm with Alpha-Beta Pruning}
The core of the agent is a recursive Minimax algorithm. To reduce the number of nodes explored, we implemented Alpha-Beta pruning. This optimization allows the search to skip branches that are guaranteed to be worse than previously explored options. The agent searches to a fixed depth of \texttt{MAX\_DEPTH = 12}, which is sufficient to cover the entire relevant game tree for a $2 \times 2$ Mancuna board.

\subsection{Transposition Table (Memoization)}
To handle the possibility of reaching the same board state through different move sequences (transpositions), we implemented a transposition table using a \texttt{Dictionary<(int Hash, int Turn), (double Value, int Depth, HashFlag Flag)>}.

\begin{lstlisting}[language={[Sharp]C}, caption=Transposition Table Lookup Implementation]
var key = (board.hash(), 0); // 0 for Max, 1 for Min

if (memo.TryGetValue(key, out var cached))
{
    if (cached.Depth >= remainingDepth)
    {
        if (cached.Flag == HashFlag.Exact)
            return cached.Value;
        if (cached.Flag == HashFlag.LowerBound && cached.Value >= beta)
            return cached.Value;
        if (cached.Flag == HashFlag.UpperBound && cached.Value <= alpha)
            return cached.Value;
    }
}
\end{lstlisting}

\begin{itemize}
    \item \textbf{Hashing:} The table uses the board's internal hash function combined with a turn indicator (0 for the maximising player, 1 for the minimising player), ensuring that the same board state reached at different turns is stored separately.
    \item \textbf{Hash Flags:} We store three types of entries:
    \begin{itemize}
        \item \textbf{Exact:} The true minimax value of the node.
        \item \textbf{LowerBound:} A value resulting from a Beta-cutoff; the true value is at least this large.
        \item \textbf{UpperBound:} A value resulting from an Alpha-cutoff; the true value is at most this large.
    \end{itemize}
\end{itemize}

\subsubsection{Excluding Terminal Scores from the Cache}
A critical detail of our transposition table is that \emph{winning and losing scores are never cached}. The guard condition \texttt{if (v < 9000 \&\& v > -9000)} ensures that only heuristic (non-terminal) values are stored.

The reason is that terminal scores are \textit{ply-dependent}: a win discovered at ply~5 is scored as $10000 - 5 = 9995$, whereas the same winning position reached at ply~3 in a later search yields $10000 - 3 = 9997$. If we cached the ply-5 score and later retrieved it at ply~3, the agent would behave as though the win is two moves further away than it actually is. This biases move ordering: the agent might choose a slower path to victory not because it is strategically better, but because the cached score artificially made the faster path look worse. By never caching terminal scores, every branch that reaches a win or loss evaluates it fresh with the correct ply distance, preserving the guarantee that the agent always chooses the fastest available win (or the slowest available loss).

\subsection{Evaluation Function and the Horizon Effect}
\label{sec:eval}
When the search reaches a leaf node (either a terminal state or the maximum search depth), it evaluates the board using the following heuristic:

\[ \text{Score} = (\text{Self Score} - \text{Opponent Score}) \times 10 \]

For terminal states the function distinguishes between wins, losses, and ties by incorporating the \textit{ply} (the number of half-moves from the root):
\begin{itemize}
    \item \textbf{Win:} $10000 - \text{ply}$
    \item \textbf{Loss:} $-10000 + \text{ply}$
    \item \textbf{Tie:} $0$
\end{itemize}

\subsubsection{Addressing the Horizon Effect}
The horizon effect is a well-known failure mode of depth-limited search: when a negative outcome lies just beyond the search horizon, the agent may take a sequence of neutral or sacrificial moves that push the inevitable loss one ply further out each time, appearing to ``avoid'' it without actually doing so.

Our evaluation function counters this in two ways:

\begin{enumerate}
    \item \textbf{Faster wins are preferred.} A win at ply~3 scores $10000 - 3 = 9997$, while a win at ply~5 scores $10000 - 5 = 9995$. The maximising player will always prefer the path with the higher score, which is the one that wins soonest.

    \item \textbf{Slower losses are preferred.} A loss at ply~5 scores $-10000 + 5 = -9995$, while a loss at ply~3 scores $-10000 + 3 = -9997$. The maximising player will always prefer the less negative score, meaning it will choose moves that delay the loss as long as possible---but crucially, it will not \emph{manufacture} a delay by making moves that are otherwise harmful, because any move that lets the opponent gain tokens worsens the heuristic score at non-terminal nodes. The agent thus has no incentive to stall in a position that is already lost; it simply maximises the number of moves before defeat.
\end{enumerate}

The gap between the maximum heuristic value ($7 \times 10 = 70$, since the board has at most 8 tokens) and the minimum terminal value ($\pm 9000$) is large enough that no sequence of heuristic scores can ever exceed a terminal score, ensuring terminal outcomes always dominate the search.

\section{Results}
To evaluate the performance of the \texttt{PlayerMancuna} agent, we conducted an extensive experiment consisting of 20,000 matches against an \texttt{Expert} player. The results demonstrate a high level of stability and competitiveness.

\begin{table}[h]
    \centering
    \begin{tabular}{|l|c|c|c|}
        \hline
        \textbf{Player} & \textbf{Wins} & \textbf{Total Score} & \textbf{Total Moves} \\
        \hline
        MancunaPlayer0 (Our Agent) & 0 & 100,000 & 250,000 \\
        Expert Player & 0 & 60,000 & 260,000 \\
        \hline
        \textbf{Ties} & \multicolumn{3}{c|}{\textbf{20,000}} \\
        \hline
    \end{tabular}
    \caption{Experimental results after 20,000 matches.}
    \label{tab:results}
\end{table}

The experiment shows a 100\% tie rate over the 20,000 matches. Given the limited state space of the $2 \times 2$ Mancuna board with only 8 tokens, these results suggest that both agents are playing optimally or near-optimally, consistently reaching a stalemate or a repeating board state that triggers a tie. Our agent maintained a higher cumulative score compared to the Expert player, indicating a strong defensive and offensive balance.

\section{Conclusion}
The implementation of the \texttt{PlayerMancuna} agent successfully combines classic game search algorithms with modern optimizations. By using Alpha-Beta pruning and a transposition table, the agent can search efficiently and avoid redundant computation for repeated board states. The evaluation function's ply-adjusted terminal scores address the horizon effect, ensuring the agent always prefers faster wins and delays losses without being deceived into artificial stalling. The deliberate exclusion of terminal scores from the transposition table preserves correct ply-distance reasoning, preventing the cache from biasing the agent toward suboptimal paths.

\end{document}
